%------------------------------------------------------------------------------%
% Luis A. D'Afonseca
%
%------------------------------------------------------------------------------%
\documentclass[a4paper,12pt,fleqn]{article}

\usepackage{style_avaliacao}
\usepackage{systeme}

\ShowAnswers

%------------------------------------------------------------------------------%
\begin{document}

\makeHeader{Integrais e Séries}{Prova 3 -- Turma 5}{09/11/2023}

% 01
%------------------------------------------------------------------------------%
\exercise{25}
Considerando a série
\[
  \sum_{n=1}^{\infty} \frac{5}{n^2 + 5n + 6}
\]
\begin{enumerate}[label=\alph*)]
  \item Calcule a soma da série
  \item Calcule o valor do erro ao aproximar a soma por \(S_{100}\)
\end{enumerate}

\begin{answer}
  \begin{enumerate}[label=\alph*)]
    \item Vamos verificar que a série é telescópica.

    Calculando as raízes de \(x^2 + 5x + 6=0\), temos
    \[
      \Delta
      = b^2 - 4ac
      = 5^2 - 4\times 1\times 6
      = 25 - 24
      = 1
    \]
    \[
      x
      = \frac{-b\pm \sqrt{\Delta}}{2a}
      = \frac{-5\pm 1}{2}
      \qquad
      x_1 = \frac{-6}{2} = -3
      \qquad
      x_2 = \frac{-4}{2} = -2
    \]
    Portanto
    \[
      x^2 + 5x + 6 = (x+2)(x+3)
    \]
    e podemos escrever o termo geral da série como
    \[
      a_n
      = \frac{5}{n^2 + 5n + 6}
      = \frac{5}{(n+2)(n+3)}
    \]
    Usando frações parciais
    \[
      \frac{5}{(n+2)(n+3)}
      = \frac{A}{n+2} + \frac{B}{n+3}
      = \frac{A(n+3) + B(n+2)}{(n+2)(n+3)}
    \]
    Portanto
    \[
      5
      = A(n+3) + B(n+2)
      = An + 3A + Bn + 2B
      = (A+B)n + (3A+2B)
    \]
    \[
      \systeme[AB]
      {
        A+B=0,
        3A+2B=5
      }
    \]
    \[
      B = -A
      \qquad\qquad
      3A-2A = 5
      \qquad\qquad
      A = 5
      \qquad\qquad
      B = -5
    \]
    Assim
    \[
      a_n = \frac{5}{n+2} - \frac{5}{n+3}
    \]
    O que comprova que a série é telescópica. Consequentemente, as somas parciais são
    \[
      S_n = \frac{5}{3} - \frac{5}{n+3}
    \]
    Calculando o limite
    \[
      S
      = \lim_{n\to\infty} S_n
      = \lim_{n\to\infty} \frac{5}{3} - \frac{5}{n+3}
      = \frac{5}{3}
    \]
    \item
    \[
      R_{100}
      = S - S_{100}
      = \frac{5}{3} - \left(\frac{5}{3} - \frac{5}{100+3}\right)
      = \frac{5}{103}
    \]
  \end{enumerate}
\end{answer}

% 02
%------------------------------------------------------------------------------%
\exercise{25}
Calcule o limite da sequência
\(
a_n = \begin{cases}
  \frac{n}{\pi-n^2}, & \; n \text{ par}   \\[2mm]
  \;\;\; 2^{-n},    & \; n \text{ ímpar}
\end{cases}
\)

\begin{answer}
  Observando que
  \(\frac{n}{\pi-n^2}<0\) para todo \(n>1\) e
  \(2^{-n}\) é sempre positivo, temos que, para \(n>1\),
  \[
    \frac{n}{\pi-n^2} \;\leq\; a_n \;\leq\; 2^{-n}
  \]
  Como
  \[
    \lim_{n\to\infty} \frac{n}{\pi-n^2}
    = \lim_{n\to\infty} \frac{\sfrac{1}{n}}{\sfrac{\pi}{n}-1}
    = 0
  \]
  e
  \[
    \lim_{n\to\infty} 2^{-n}
    = \lim_{n\to\infty} \left(\frac{1}{2}\right)^n
    = 0
  \]
  o \textbf{Teorema do Confronto} garante que \(\lim_{n\to 0} a_n = 0\)
\end{answer}

% 03
%------------------------------------------------------------------------------%
\exercise{25}
Verifique se a série converge
\[
  \sum_{n=1}^\infty \frac{\sin^2(n)}{n^\pi}
\]

\begin{answer}
  Como os termos são não-negativos podemos usar o teste da comparação.
  Observando que
  \(
    \sin^2(n) \leq 1
  \)
  temos que
  \[
    a_n = \frac{\sin^2(n)}{n^\pi} \leq \frac{1}{n^\pi}
  \]
  Como \(b_n = \frac{1}{n^\pi} \) são os termos de uma $p$-série convergente
  (\(p=\pi>1\))
  o teste da comparação garante que a série
  \[
    \sum_{n=1}^\infty \frac{\sin^2(n)}{n^\pi}
  \]
  converge.
\end{answer}

% 04
%------------------------------------------------------------------------------%
\exercise{25}
Use o teste da integral para verificar se a soma da série existe
\[
  \sum_{n=1}^{\infty} n e^{-n^2}
\]
Não se esqueça de verificar que as condições do teste são satisfeitas.

\begin{answer}
  Vamos usar a função real
  \[
    f(x) = x e^{-x^2}
  \]

  Precisamos verificar as três condições do teste
  \begin{enumerate}[label=\roman*)]
    \item \(f(x)>0\) em \([1, \infty)\) \\[2mm]
      Como a exponencial é sempre positiva e \(x\geq 1\) temos que \(f(x)>0\)

  \item \(f(x)\) é contínua em \([1, \infty)\) \\[2mm]
    Como a exponencial é contínua, \(f(x)\) é contínua em \(\R\)

  \item \(f(x)\) é decrescente em  \([1, \infty)\) \\[2mm]
    A derivada de \(f(x)\) é
    \begin{align*}
    f'(x)
    & = (x)'e^{-x^2} + x \left(e^{-x^2}\right)' \\
    & = e^{-x^2} + x e^{-x^2} \left(-x^2\right)' \\
    & = e^{-x^2} + x e^{-x^2} \left(-2x\right) \\
    & = e^{-x^2} -2 x^2 e^{-x^2} \\
    & = e^{-x^2}\left(1 -2 x^2\right)
    \end{align*}
    portanto \(f'(x)<0\) para \(x>1\), assim \(f\) é decrescente no intervalo \([1, \infty)\)
  \end{enumerate}
  Calculando agora a primitiva de \(f\)
  \[
    F(x) = \int x e^{-x^2} dx
  \]
  Fazendo a mudança de variáveis
  \[
    u = -x^2
    \qquad
    du = -2xdx
    \qquad
    dx = \frac{du}{-2x}
  \]
  \[
    F(x)
    = \int x e^{-x^2} dx
    = \int x e^{u} \frac{du}{-2x}
    = \frac{-1}{2} \int e^{u} du
    = \frac{-e^u}{2} + c
    = \frac{-e^{-x^2}}{2} + c
  \]
  Calculando a integral imprópria
  \[
    A
    = \int_1^\infty f(x) dx
    = \lim_{b\to\infty}\int_1^b f(x)dx
    = \lim_{b\to\infty} F(x)\evalat{1}{b}
    = \lim_{b\to\infty} \left(\frac{-e^{-b^2}}{2} - \frac{-e^{-1^2}}{2}\right)
    = \frac{e^{-1}}{2}
    = \frac{1}{2e}
  \]
  Como a integral converge a série também converge e portanto sua soma existe.
\end{answer}

%------------------------------------------------------------------------------%
\end{document}
%------------------------------------------------------------------------------%
